\documentclass[a4paper,11pt]{article}

\usepackage[utf8]{inputenc}
\usepackage[T1]{fontenc}
\usepackage{lmodern}
\usepackage[spanish,es-nodecimaldot,es-noshorthands]{babel}
\usepackage{amsmath, amssymb, mathtools}
\usepackage[left=2.3cm,right=2.3cm,top=2.3cm,bottom=2.3cm]{geometry}
\usepackage{microtype}
\usepackage{parskip}
\usepackage{xcolor}
\usepackage{booktabs}
\usepackage{tabularx}
\usepackage{array}
\usepackage{enumitem}
\usepackage[most]{tcolorbox}
\usepackage{tikz}
\usetikzlibrary{arrows.meta,positioning}
\usepackage{pgfplots}
\pgfplotsset{compat=1.18}
\usepackage{hyperref}
\hypersetup{colorlinks=true, linkcolor=blue!60!black, urlcolor=blue!60!black, citecolor=green!45!black}

\addto\captionsspanish{\renewcommand{\tablename}{Tabla}}
\setlist{itemsep=2pt, topsep=3pt}
\renewcommand{\arraystretch}{1.18}
\newcolumntype{Y}{>{\raggedright\arraybackslash}X}

\newtcolorbox{idea}[1][Idea clave]{colback=blue!4, colframe=blue!45!black, fonttitle=\bfseries, title=#1, boxrule=0.6pt, arc=2pt}
\newtcolorbox{ojo}[1][Ojo]{colback=orange!6, colframe=orange!70!black, fonttitle=\bfseries, title=#1, boxrule=0.6pt, arc=2pt}
\newtcolorbox{ejemplo}[1][Ejemplo]{colback=green!4, colframe=green!40!black, fonttitle=\bfseries, title=#1, boxrule=0.6pt, arc=2pt}
\newcommand{\rec}{\ensuremath{\star}}
\newcommand{\pos}{\ensuremath{\mathcal{P}}}
\newcommand{\cel}[3]{\fill[#3,draw=white,thin] (#1,#2) rectangle ++(1,1);}
\newcommand{\ctr}{blue!35}\newcommand{\cva}{orange!60}\newcommand{\cte}{red!55}\newcommand{\ccx}{gray!35}

\title{\textbf{Partir los datos y describir las películas}\\[4pt]
\large Guía para responder a los bloques D (split) y E (features de la torre de películas) del recomendador de CPIA}
\author{Arnau Montull \quad para \quad Uriel Martínez}
\date{16 de octubre de 2026}

\begin{document}
\maketitle

\begin{abstract}
\noindent Esta guía reúne lo necesario para decidir dos cosas: \textbf{(D)} cómo partimos los ratings en entrenamiento,
validación y test sin hacernos trampas, y \textbf{(E)} qué información sobre cada película (sinopsis, géneros, reparto,
año\dots) recibe la torre de ítem. Cada concepto viene con un ejemplo y con cifras medidas en la submuestra de desarrollo
(3\,154\,155 ratings, 20\,085 usuarios, 22\,787 películas con rating): cuántos usuarios y películas de test no aparecen en
entrenamiento con cada tipo de split, cuánta fuga habría, qué cobertura tiene cada feature y qué pasa con las películas
de estreno reciente. Al final están las preguntas D1--D10 y E1--E9 con sus opciones, implicaciones (modelo, pérdida,
evaluación, entrega E2), qué decisiones del bloque C les afectan y mi recomendación (\rec), que \emph{no} es una decisión.
Las cifras salen de \texttt{python -m src.analisis.split\_y\_features} (\texttt{reports/split\_y\_features.md}); la versión
corta, con la hoja de respuestas, está en \texttt{docs/planning\_bloque\_D\_E.md}.
\end{abstract}

\tableofcontents
\newpage

% =====================================================================
\section{Dónde estamos y qué se decide aquí}
% =====================================================================

\subsection{Lo ya decidido (bloques A, B y C)}
\begin{itemize}
  \item \textbf{A (datos)}: solo \texttt{ml-32m/ratings.csv} (1995--2023); de Kaggle/TMDB solo metadatos. Submuestra de
        desarrollo: el 10\,\% de los usuarios enteros, estratificada; el conjunto completo, para los resultados finales.
  \item \textbf{B (limpieza)}: fuera las películas con menos de 20 ratings (se pierde el 0,86\,\% de los ratings);
        quedan 31\,725\,154 ratings, 200\,948 usuarios y 23\,335 películas. Duplicados (mismo \texttt{tmdbId})
        fusionados. Sinopsis para el 99,5\,\% de las películas. \textbf{B3 sigue pendiente} (géneros de MovieLens o
        TMDB): lo tratamos en E5.
  \item \textbf{C (señal objetivo)}, provisional salvo lo marcado:
\end{itemize}
\begin{center}
\small
\begin{tabularx}{\linewidth}{lY}
\toprule
Pregunta & Decisión \\
\midrule
C1 b & Implícito binario; pérdida contrastiva con negativos del lote y corrección logQ. \\
C2 & Positivo si $r\ge\min(T;\bar r_u)$ con \textbf{$T=4{,}0$ (confirmado)}. \\
C3 c & Anti-preferencias en la torre de usuario; «nota baja» $=r\le\bar r_u-1$. \\
C4 d & Peso $w=1+\max(0,\,r-\bar r_u)$ multiplicando la pérdida. \\
C5 a & Cada rating positivo es un ejemplo, sin tope por usuario. \\
C6 a & $\bar r_u$ cruda, calculada solo con train. \\
C7 c & Relevante en test: con la definición de train \emph{y} con $r\ge4{,}0$. \\
C8 b & Usuarios sin positivos en train o test se quedan en train y salen de las métricas. \\
C9 c & Peso $\times0{,}75$ a los ratings de las primeras 24\,h de cada usuario. \\
C10 a, C11 b, C12 a & Nada para la escala; quitar ratings anteriores al estreno; re-ratings: la última. \\
C13 b & $P(i)=$ positivos de $i$ en train $/$ total, con $+1$ de suavizado. \\
C14 c & Enmascarar en la softmax todo lo que el usuario ha valorado en train. \\
\bottomrule
\end{tabularx}
\end{center}

\subsection{La arquitectura que prometimos en E1}
\begin{center}
\begin{tikzpicture}[node distance=5mm and 9mm, every node/.style={font=\small},
  caja/.style={draw, rounded corners=2pt, align=center, minimum height=9mm, fill=#1},
  caja/.default=white, flecha/.style={-{Stealth[length=2.2mm]}, thick}]
  \node[caja=blue!8] (hist) {Historial del usuario\\(notas)};
  \node[caja=blue!18, right=of hist] (tu) {Torre de usuario\\$\Psi_U$};
  \node[caja=green!8, below=12mm of hist] (meta) {Sinopsis, géneros,\\director, año};
  \node[caja=green!18, right=of meta] (ti) {Torre de ítem\\$\Psi_I$};
  \node[caja=orange!12, right=14mm of tu, yshift=-10.5mm] (dot) {$\langle \mathbf{u}, \mathbf{v}_i\rangle$\\+ ANN (HNSW)};
  \node[caja=orange!22, right=of dot] (mmr) {200 candidatos\\$\rightarrow$ MMR $\rightarrow$ 20};
  \draw[flecha] (hist) -- (tu); \draw[flecha] (meta) -- (ti);
  \draw[flecha] (tu) -- node[above, sloped]{$\mathbf{u}$} (dot); \draw[flecha] (ti) -- node[below, sloped]{$\mathbf{v}_i$} (dot);
  \draw[flecha] (dot) -- (mmr);
  \node[below=1mm of meta, font=\footnotesize\itshape, text=red!60!black] {aquí entra el bloque E};
  \node[above=1mm of hist, font=\footnotesize\itshape, text=red!60!black] {el bloque D decide qué historial y qué test};
\end{tikzpicture}
\end{center}

En E1 prometimos: Two-Tower con pérdida contrastiva y corrección logQ \cite{yi2019}, búsqueda aproximada (HNSW) para
sacar 200 candidatos, MMR \cite{carbonell1998} para dejar 20, embedding de la sinopsis con MiniLM ($d=384$), géneros y
director en multi-hot, y el \emph{export} de Letterboxd de un usuario nuevo como caso real.

\subsection{Qué decide cada bloque, y el supuesto que he tomado para E}
\begin{itemize}
  \item \textbf{Bloque D = split train/validación/test}: cómo se reparten los ratings, qué usuarios y películas ve el
        modelo al entrenar, con qué se valida, cómo se mide (métricas y $K$) y cómo se hace reproducible.
  \item \textbf{Bloque E = features de la torre de películas}. Esto no aparece definido en la entrega E1 ni en el PDF
        de etapas del proyecto de la asignatura; el supuesto sale de \texttt{docs/\allowbreak planning\_\allowbreak bloques\_\allowbreak D\_H.md} (bloque E:
        embedding de sinopsis, créditos, keywords y numéricas). El bloque F (torre de usuario) queda fuera: aquí solo se
        decide qué sabe el modelo de \emph{cada película}.
\end{itemize}

\begin{ojo}[Tres cosas que se cruzan]
\begin{itemize}[leftmargin=*]
  \item \textbf{D1 y F1} van juntos: un split por usuarios solo tiene sentido si la torre de usuario lee el historial y
        no un identificador de usuario. Y D1 define el escenario de Letterboxd (un usuario que no estaba en el
        entrenamiento).
  \item \textbf{D8 y E6} van juntos: el arranque en frío de películas solo se puede \emph{medir} si el split deja
        películas sin entrenar; y solo se puede \emph{resolver} si la torre de ítem funciona con solo contenido.
  \item \textbf{D3 y E1--E4}: todo lo que se ajusta (vocabularios, frecuencias mínimas, medias para normalizar) tiene
        que respetar el patrón \emph{fit/transform}.
\end{itemize}
\end{ojo}

% =====================================================================
\section{Bloque D: cómo partir los datos}\label{sec:D}
% =====================================================================

\subsection{Para qué sirve cada partición}
\begin{itemize}
  \item \textbf{Train}: con él se ajustan los pesos de las torres \emph{y} todo lo que se «calcula de los datos»
        ($\bar r_u$, $P(i)$ de logQ, la máscara de C14, vocabularios\dots).
  \item \textbf{Validación}: sirve para elegir hiperparámetros (temperatura, tamaño, tasa de aprendizaje\dots) y parar
        el entrenamiento a tiempo (\emph{early stopping}). Cada vez que miramos su resultado y decidimos algo, «se
        gasta» un poco.
  \item \textbf{Test}: se mira una vez, al final, para decir cuánto vale el modelo. Si se usa para decidir, deja de
        ser test.
\end{itemize}

\begin{idea}
Un split es una \emph{hipótesis sobre el futuro}: dice qué información estará disponible cuando haya que recomendar y
qué no. Cuanto más se parezca el split al uso real (un usuario nuevo trae su historial y pide recomendaciones), más
honestas son las métricas. En esta guía, ese uso real es el de Letterboxd.
\end{idea}

\subsection{Los cuatro tipos de split}
\begin{center}
\begin{tikzpicture}[x=0.26cm, y=0.3cm, font=\footnotesize]
  \tikzset{tr/.style={fill=blue!35, draw=white, thin}, va/.style={fill=orange!60, draw=white, thin},
           te/.style={fill=red!55, draw=white, thin}, cx/.style={fill=gray!35, draw=white, thin},
           vacio/.style={fill=white, draw=gray!30, thin}}
  % (a) aleatorio
  \begin{scope}
    \pgfmathsetseed{11}
    \foreach \y in {0,...,5}{\foreach \x in {0,...,11}{
      \pgfmathparse{rnd}\let\r\pgfmathresult
      \ifdim\r pt<0.8pt \cel{\x}{\y}{\ctr}\else\ifdim\r pt<0.9pt \cel{\x}{\y}{\cva}\else\cel{\x}{\y}{\cte}\fi\fi}}
    \node at (6,-2.2) {(a) aleatorio};
  \end{scope}
  % (b) temporal global
  \begin{scope}[xshift=3.8cm]
    \foreach \y/\ini in {0/0,1/0,2/2,3/5,4/7,5/9}{\foreach \x in {\ini,...,11}{
      \ifnum\x<9 \cel{\x}{\y}{\ctr}\else\ifnum\x<10 \cel{\x}{\y}{\cva}\else\cel{\x}{\y}{\cte}\fi\fi}}
    \node at (6,-2.2) {(b) temporal global};
  \end{scope}
  % (c) último por usuario
  \begin{scope}[xshift=7.6cm]
    \foreach \y/\fin in {0/11,1/9,2/11,3/7,4/11,5/10}{\foreach \x in {0,...,\fin}{
      \ifnum\x<\fin \ifnum\x<\numexpr\fin-1\relax \cel{\x}{\y}{\ctr}\else\cel{\x}{\y}{\cva}\fi\else\cel{\x}{\y}{\cte}\fi}}
    \node at (6,-2.2) {(c) último de cada usuario};
  \end{scope}
  % (d) por usuarios
  \begin{scope}[xshift=11.4cm]
    \foreach \y in {2,...,5}{\foreach \x in {0,...,11}{\cel{\x}{\y}{\ctr}}}
    \foreach \y/\g in {0/\cva,1/\cte}{\foreach \x in {0,...,11}{\ifnum\x<9 \cel{\x}{\y}{\ccx}\else\cel{\x}{\y}{\g}\fi}}
    \node at (6,-2.2) {(d) por usuarios};
  \end{scope}
  % leyenda
  \begin{scope}[yshift=-1.9cm, xshift=0cm]
    \fill[blue!35] (0,0) rectangle ++(1.2,1.2); \node[right] at (1.4,0.6) {train};
    \fill[orange!60] (10,0) rectangle ++(1.2,1.2); \node[right] at (11.4,0.6) {validación};
    \fill[red!55] (24,0) rectangle ++(1.2,1.2); \node[right] at (25.4,0.6) {test};
    \fill[gray!35] (33,0) rectangle ++(1.2,1.2); \node[right] at (34.4,0.6) {contexto (entrada del modelo, no se predice)};
  \end{scope}
  \node[rotate=90] at (-1.6,3) {usuarios};
  \draw[-{Stealth}] (0,-0.9) -- (12,-0.9) node[right, font=\tiny] {tiempo};
\end{tikzpicture}
\end{center}

\begin{itemize}
  \item \textbf{(a) Aleatorio por rating}: cada rating cae al azar en una partición. Fácil, pero entrenamos con ratings
        posteriores al que predecimos. Es el tipo que \cite{covington2016} descarta porque «filtra información del futuro».
  \item \textbf{(b) Temporal global}: un corte por fecha; lo anterior a train, lo posterior a test. No hay fuga
        temporal, pero casi todo el test pertenece a usuarios que se registraron después del corte.
  \item \textbf{(c) Último de cada usuario} (\emph{leave-last-out}): su última interacción a test, la anterior a
        validación. Es el protocolo de NCF y SASRec \cite{he2017,kang2018}.
  \item \textbf{(d) Por usuarios} (generalización fuerte): un grupo de usuarios entero se aparta del entrenamiento; de
        cada uno se da al modelo el 80\,\% de su historial (\emph{fold-in}) y se predice el 20\,\% restante. Es el
        protocolo de Mult-VAE y de los trabajos que le siguen \cite{liang2018,steck2019}, y es exactamente el
        escenario de un usuario nuevo de Letterboxd.
\end{itemize}

\begin{ejemplo}[Ejemplo: Ana con cada split]
Ana tiene 10 ratings. En (a) el 8 de ellos va a train, uno a val y uno a test \emph{al azar}: su rating más reciente
puede estar en train, y uno antiguo en test. En (c) los 8 primeros van a train, el 9.º a val y el 10.º a test. En (b) los
que haya antes del corte van a train y los que haya después, a test (si Ana se registró después, \emph{todo} su historial
cae en test y el modelo no la ha visto nunca). En (d), o Ana entera es de train, o Ana es de test: los 8 primeros ratings
entran en la torre de usuario y el modelo tiene que acertar los 2 últimos.
\end{ejemplo}

\subsection{Lo que dicen los datos sobre cada split}
Todos los splits son 80/10/10 sobre la submuestra de desarrollo, con la semilla 42. «Positivo» es el de C2
($r\ge\min(4{,}0;\bar r_u)$, con $\bar r_u$ de la historia visible).
\begin{center}
\footnotesize
\begin{tabular}{lrrrr}
\toprule
 & Aleatorio & Último/usuario & Temporal global & Por usuarios \\
\midrule
Ratings en train & 2\,523\,007 & 2\,505\,121 & 2\,523\,324 & 2\,528\,833 \\
Ratings en val & 315\,963 & 324\,517 & 315\,429 & 64\,835 \\
Ratings en test & 315\,185 & 324\,517 & 315\,402 & 61\,867 \\
Usuarios con ratings en test & 19\,690 & 20\,085 & 2\,259 & 2\,069 \\
\addlinespace
Usuarios de test sin nada en train & 0 & 0 & 1\,809 & 2\,069 \\
\% ratings de test de usuarios sin train & 0 & 0 & 86,3 & 100 \\
\% ratings de test en películas sin train & 0,06 & 0,43 & 13,24 & 0,15 \\
\% ratings de test en películas con $<5$ en train & 0,90 & 2,73 & 15,63 & 1,54 \\
\% ratings de test en películas con $<20$ en train & 3,32 & 7,72 & 19,43 & 5,35 \\
\addlinespace
\% de test \emph{anterior} al último rating visible & 98,9 & 0 & 0 & 0 \\
\% de test el mismo día del último rating visible & 1,0 & 50,2 & 0 & 62,2 \\
\bottomrule
\end{tabular}
\end{center}
Corte del temporal global: train hasta el 22/07/2018, validación hasta el 22/10/2020, test después. (En el conjunto
completo los cortes son el 01/09/2018 y el 21/10/2020: casi los mismos.)

\begin{center}
\begin{tikzpicture}
\begin{axis}[ybar, width=12cm, height=5.4cm, bar width=9pt, ymin=0, ymax=100, enlarge x limits=0.18,
  symbolic x coords={Aleatorio, Último por usuario, Temporal global, Por usuarios}, xtick=data,
  ylabel={\% de los ratings de test}, ylabel style={font=\small}, tick label style={font=\footnotesize},
  legend style={at={(0.02,0.97)}, anchor=north west, font=\footnotesize}, nodes near coords, nodes near coords style={font=\tiny}]
\addplot[fill=blue!35] coordinates {(Aleatorio,0)(Último por usuario,0)(Temporal global,86.3)(Por usuarios,100)};
\addplot[fill=orange!65] coordinates {(Aleatorio,0.9)(Último por usuario,2.7)(Temporal global,15.6)(Por usuarios,1.5)};
\legend{De usuarios sin nada en train, En películas con $<5$ ratings en train}
\end{axis}
\end{tikzpicture}
\end{center}

Lo que enseña:
\begin{itemize}
  \item \textbf{Aleatorio}: el 98,9\,\% de los ratings de test son \emph{anteriores} al último rating que el modelo ve de
        ese usuario: se entrena con el futuro. Casi no hay frío, ni de usuarios ni de películas.
  \item \textbf{Temporal global}: el 86,3\,\% del test son ratings de usuarios que no están en train, y el 13,2\,\% son de
        películas sin ningún rating en train. Solo 439 de los 2\,259 usuarios de test (19,4\,\%) tienen algún positivo en
        test \emph{y} en su historia: el test mide sobre todo recomendación a usuarios nuevos sin historia. Hay 354 usuarios
        repartidos entre las tres particiones: son los que llevan años en el sistema.
  \item \textbf{Último por usuario}: realista en el orden, pero el 50,2\,\% del test se hizo el mismo día que el último
        rating de train. Y el corte cae en medio de un empate de segundo en el 24,4\,\% de los usuarios (\S\ref{sec:primerdia}).
  \item \textbf{Por usuarios}: el 100\,\% de los usuarios de test son nuevos para el modelo (como en Letterboxd), solo
        0,15\,\% de los ratings de test son de películas sin entrenamiento, y 2\,051 de los 2\,069 usuarios de test son
        elegibles para las métricas (99,1\,\%).
\end{itemize}

\subsection{El primer día: 59\,\% de usuarios sin orden real}\label{sec:primerdia}
Recordatorio de la guía C: el 62,5\,\% de los ratings se hace en las primeras 24\,h del usuario, y el 59\,\% de los
usuarios lo valora todo ese primer día. Los autores de MovieLens avisan de que las marcas de tiempo \emph{no} son la fecha
de consumo y de que los usuarios introducen muchos ratings de golpe, rellenando lo que vieron antes \cite{harper2015}.
Consecuencias para el split:
\begin{itemize}
  \item «Lo siguiente que ve el usuario» no es lo siguiente que \emph{vio}: es lo siguiente que \emph{puntuó}. Para el
        59\,\% de usuarios, el orden interno es el de teclear la lista de recuerdos.
  \item En el split «último por usuario» (10\,\%+10\,\%) el corte cae dentro de un empate de segundo en el
        \textbf{17,7\,\%} de los usuarios (train/val) y en el \textbf{16,3\,\%} (val/test): en total, el 24,4\,\% de los
        usuarios tiene al menos un corte «arbitrario» (desempata el \texttt{movieId}).
  \item El primer día pesa menos en test que en train (36,7\,\% de los ratings de test frente al 68,4\,\% de los de train
        en el split por último 10\,\%): se evalúa más con lo posterior, y eso es lo que queremos. Pero con C9 c (peso $\times0{,}75$
        a las primeras 24\,h) el entrenamiento ya \emph{baja} el peso de esos ratings: la evaluación debe ir
        \emph{sin ponderar}.
\end{itemize}

\begin{ojo}[¿«Cronológico» es más honesto?]
Para los usuarios de varias sesiones, sí. Para el 59\,\% que lo hace todo de golpe, predecir «el final de la lista de
recuerdos» no es predecir el futuro del gusto; es predecir un orden de tecleo. No es una trampa, pero sí es ruido: por eso
conviene decir en el informe que el orden dentro del primer día no es información, y repetir el resultado con un split
aleatorio dentro del usuario como control (pregunta D2).
\end{ojo}

\subsection{Fuga silenciosa: lo que se calcula con todos los datos}
La fuga más peligrosa no es mezclar ratings, sino calcular estadísticas con test sin darse cuenta. En nuestra
canalización hay tres: la media del usuario (C6), la frecuencia $P(i)$ de logQ (C13) y la máscara de C14. Con el split
«último por usuario» medimos qué pasaría si se calcularan con todo:
\begin{itemize}
  \item $\bar r_u$ con todo el historial frente a solo train: diferencia absoluta media de 0,062 (p95 0,178, máx. 0,92).
  \item El \textbf{1,46\,\%} de los ratings de train (36\,534) cambiaría de etiqueta positivo/no positivo, y el 1,25\,\% de
        los de test.
  \item 1\,730 películas tienen algún positivo en la muestra completa pero ninguno en train. Con $P(i)$ calculada con todo
        se les daría un $\log P(i)$ finito que solo existe porque hay positivos en test (la correlación de Spearman entre
        ambos conteos es 0,972: no es un desastre, es un sesgo pequeño pero injustificable).
  \item La máscara de C14 (lo valorado \emph{en train}) no puede incluir lo valorado en test: sería quitar de la softmax
        justo las películas que hay que acertar.
\end{itemize}
La regla: ajustar con train y transformar el resto (patrón \emph{fit/transform}). Vale también para vocabularios
(directores, actores), la frecuencia mínima de E2 y las medias y desviaciones con las que se normalizan año y duración (E4).

\subsection{Usuarios y películas que no están en train}
\textbf{Usuarios}: en el split por usuarios, todos los de test son nuevos por construcción; el modelo solo los conoce por su historial
(contexto). Eso exige una torre de usuario que lea el historial (bloque F). En el temporal global, el 86,3\,\% del test es de
usuarios sin historia visible, y no se les puede evaluar con ninguna métrica de recomendación personalizada.

\textbf{Películas}: en nuestros splits con usuarios en train, apenas el 0,15\,\%--0,43\,\% de los ratings de test son de películas
que train no ha visto nunca (13,2\,\% en el temporal global). Con eso, \emph{no podemos medir el arranque en frío de películas}
(la promesa de E1) con un split de ratings: tendríamos que retirar películas adrede (pregunta D8). Y ojo con la submuestra: el
filtro de 20 ratings se aplicó al conjunto completo; en la submuestra, 15\,422 películas de las 23\,335 tienen menos de 20
ratings en train, y 4\,354 no tienen ningún positivo ($P(i)$ solo existe por el suavizado $+1$). En el conjunto completo, con
diez veces más usuarios, habrá menos películas en esa situación: el frío de la submuestra es peor que el real.

\subsection{Métricas sobre el catálogo completo y qué $K$ usar}
Las métricas son las de la guía C: Recall@$K$ y NDCG@$K$ \cite{liang2018} (esta última, premiando acertar arriba), sobre
\emph{todo} el catálogo (22\,787 películas en la submuestra, quitando lo que el usuario ya tiene en su historial visible),
no sobre 100 negativos al azar. Esa práctica de muestrear 100 negativos es la que usan SASRec y NCF por coste
\cite{kang2018,he2017}, pero aquí el catálogo cabe: 2\,056 usuarios de validación $\times$ 22\,787 películas son 47 millones
de productos escalares. Además de las dos de ranking, se miden \textbf{cobertura@$K$} (qué fracción del catálogo llega a
aparecer en alguna lista) y \textbf{novedad@$K$} (por ejemplo, $-\log_2 P_{\text{train}}(i)$ media de lo recomendado).

\textbf{¿Qué $K$?} El sistema tiene dos etapas: HNSW devuelve 200 candidatos y MMR deja 20. Cada etapa tiene su $K$:
Recall@200 mide si la primera etapa trae los relevantes (es lo que entrena la Two-Tower); NDCG@20 y Recall@20 miden la lista final
(y se calculan \emph{antes y después} del MMR). Un $K$ demasiado pequeño choca con cuántos relevantes hay: el \textbf{techo} de
Recall@$K$ es la media de $\min(K,|R_u|)/|R_u|$ (el máximo que sacaría un modelo perfecto).
\begin{center}
\small
\begin{tabular}{lrrrr}
\toprule
Techo de Recall@$K$ (\%) & Aleatorio & Último/usuario & Temporal global & Por usuarios \\
\midrule
$K=10$ & 87,5 & 88,7 & 55,1 & 77,1 \\
$K=20$ & 95,4 & 96,0 & 71,1 & 89,9 \\
$K=50$ & 99,3 & 99,4 & 86,9 & 97,8 \\
$K=200$ & 100 & 100 & 98,5 & 100 \\
\addlinespace
Positivos por usuario elegible (mediana) & 5 & 5 & 20 & 9 \\
\% de elegibles con $<3$ positivos en test & 24,9 & 20,0 & 10,3 & 6,5 \\
Usuarios elegibles & 18\,759 & 19\,396 & 439 & 2\,051 \\
\bottomrule
\end{tabular}
\end{center}
En el split por usuarios, un Recall@10 perfecto es 77,1\,\%: si se compara Recall@10 con la literatura, hay que tener presente
este techo. Con 2\,051 usuarios elegibles y 35\,019 positivos en test, el ruido por usuario es alto; hay que acompañar
cada métrica con un intervalo (\emph{bootstrap} sobre usuarios).

\subsection{Validación, tamaño de las particiones y semillas}
\begin{itemize}
  \item \textbf{Validación}: lo más coherente es una validación \emph{del mismo tipo que el test}: usuarios apartados (2\,056
        en la submuestra) con su 80\,\% de historia y el 20\,\% a predecir, para parar el entrenamiento y elegir
        hiperparámetros. Así validamos lo que vamos a usar (usuario nuevo con historial).
  \item \textbf{Tamaño}: Mult-VAE aparta 10\,000 usuarios para validación y otros 10\,000 para test de 136\,677 (el 7,3\,\%
        cada uno en ML-20M) \cite{liang2018}. Con 80/10/10 nuestro conjunto completo da $\approx$20\,000 usuarios de validación
        y otros tantos de test.
  \item \textbf{Semillas}: la semilla cambia qué usuarios caen en test. Con 5 semillas, en el split por usuarios los ratings de
        test varían $66\,565\pm1\,569$, los usuarios elegibles $2\,005\pm38$ y el \% de test en películas con $<5$ ratings en train
        $1{,}8\pm0{,}5$; en el aleatorio las cifras apenas se mueven ($315\,203\pm223$). Con tan pocos usuarios de test, un resultado
        con una sola semilla se parece a una moneda; repetir con varias semillas es parte del método.
  \item \textbf{Submuestra o completo}: la submuestra es útil para iterar; el conjunto completo da más usuarios de test
        (menos ruido) y menos frío espurio. Los cortes temporales globales de ambos coinciden casi (julio--sept. de 2018 y
        octubre de 2020).
\end{itemize}

% =====================================================================
\section{Bloque E: qué sabe la torre de ítem de cada película}\label{sec:E}
% =====================================================================

\subsection{Para qué sirve la torre de ítem}
La torre de ítem $\Psi_I$ convierte cada película en un vector $\mathbf v_i\in\mathbb R^{128}$. Los $\mathbf v_i$ no
dependen del usuario: se calculan una vez para todo el catálogo y se guardan en el índice HNSW. El contenido es lo que
permite que una película \emph{sin ratings} tenga un vector razonable: la razón para usar la arquitectura en el catálogo con cola larga
\cite{yi2019}. Yi et al. describen una torre de ítem que codifica «una gran variedad de características de contenido» y
la proponen precisamente para la escasez de datos y la distribución de potencias de los ítems.

\subsection{Cobertura de cada feature (cifras medidas)}
\begin{center}
\small
\begin{tabular}{lrr}
\toprule
Feature & \% películas & \% ratings \\
\midrule
Géneros de MovieLens & 99,0 & 99,9 \\
Géneros de TMDB & 98,7 & 99,8 \\
Géneros de alguna de las dos & 99,8 & 100,0 \\
Sinopsis & 99,5 & 99,9 \\
Director (\texttt{credits.csv}) & 84,8 & 97,0 \\
Reparto (\texttt{credits.csv}) & 83,7 & 96,9 \\
Palabras clave (\texttt{keywords.csv} o API) & 84,7 & 98,5 \\
Año de estreno & 99,7 & 99,9 \\
Duración (\texttt{runtime} de Kaggle) & 84,7 & --- \\
\bottomrule
\end{tabular}
\end{center}
«\% ratings» pondera cada película por sus ratings en la submuestra: las que más pesan en el entrenamiento están casi todas completas.

\begin{ojo}[Las películas nuevas son justo las que no tienen créditos]
\begin{center}
\small
\begin{tabular}{lrrrrr}
\toprule
Estreno & Películas & \% director & \% reparto & \% keywords & \% sinopsis \\
\midrule
$\le$1989 & 6\,336 & 98,0 & 97,1 & 87,2 & 99,7 \\
1990--99 & 3\,424 & 97,5 & 96,5 & 81,8 & 99,3 \\
2000--09 & 5\,432 & 97,8 & 96,3 & 82,5 & 99,5 \\
2010--16 & 4\,709 & 97,7 & 96,2 & 80,8 & 99,5 \\
$\ge$2017 & 3\,361 & \textbf{9,2} & \textbf{9,0} & 91,9 & 99,2 \\
\bottomrule
\end{tabular}
\end{center}
\texttt{credits.csv} es de Kaggle (2017): las películas posteriores tienen sinopsis y géneros (los pone la API de TMDB que
descargamos) pero casi nunca director ni reparto, ni duración (8,9\,\%). Si director y reparto entran como feature, el modelo ve
«desconocido» exactamente en las películas más recientes, que son las del arranque en frío. Hay que decidir qué hacer con
ello (E9).
\end{ojo}

\subsection{Sinopsis: texto a vector}
La sinopsis tiene una mediana de 46 palabras (p5: 16; p95: 114; el 0,5\,\% tiene menos de 10). Cabe sobradamente en la ventana de MiniLM
(\texttt{all-MiniLM-L6-v2}, 256 \emph{word pieces}, $\approx$190 palabras). Fuente: 19\,804 de Kaggle, 3\,391 de la API, 13 corregidas a mano y 127
sin fuente (sinopsis vacía).
\begin{itemize}
  \item \textbf{MiniLM}: red preentrenada que da un vector de 384 números por texto; se ejecuta una vez y se guarda
        (\texttt{.npy}). 23\,335 $\times$ 384 son 17,9\,MB en \texttt{float16} y 35,8\,MB en \texttt{float32}.
  \item \textbf{TF-IDF + SVD}: cuenta palabras ponderadas por su rareza y las comprime con una descomposición en
        valores singulares; no usa red neuronal y es rapidísimo, pero no sabe que «car» y «automobile» son lo mismo.
  \item Estas representaciones son \emph{del texto}, no de los ratings: ajustar el idf o la SVD con el catálogo entero no filtra
        información de test (sí lo haría un vocabulario construido con los positivos). Para el experimento de películas retiradas
        (D8), conviene ajustarlas sin esas películas, por pulcritud.
\end{itemize}
Sobre el modelo MiniLM no cito un artículo: no está en \texttt{docs/papers} y no lo he podido comprobar aquí; en la E2 habrá que
referenciar el modelo de la biblioteca \emph{sentence-transformers} con su ficha oficial.

\subsection{Géneros: la decisión B3 pendiente}
\begin{center}
\small
\begin{tabular}{lrr}
\toprule
 & MovieLens & TMDB \\
\midrule
Géneros distintos en el catálogo & 19 & 23 \\
Géneros por película (media) & 2,09 & 2,36 \\
Películas con géneros & 99,0\,\% & 98,7\,\% \\
Películas con géneros solo en esta fuente & 258 & 201 \\
\bottomrule
\end{tabular}
\end{center}
Sobre los 16 géneros comunes (Sci-Fi = Science Fiction, Children $\approx$ Family), en las 22\,838 películas que tienen las dos
listas el Jaccard medio es \textbf{0,75} y las listas son idénticas solo en el \textbf{50,2\,\%}. Hay 38 películas sin géneros
en ninguna. Solo MovieLens: Musical, Film-Noir, IMAX (un formato de proyección) y Children. Solo TMDB: Family, Music, History,
Foreign (indica que no es estadounidense) y TV Movie.

El argumento decisivo no es la cobertura: es la \textbf{inferencia}. Las películas nuevas (posteriores a 2023, y las
de Letterboxd que no están en MovieLens) no tienen géneros de MovieLens, pero sí los tienen en TMDB. Si la torre se
entrena con géneros de MovieLens, hay que reemplazarlos por otros al servir.

\subsection{Créditos y palabras clave}
\begin{itemize}
  \item \textbf{Director}: 7\,827 directores distintos; el 40,5\,\% de las películas tiene un director con 5 películas o más en el catálogo.
  \item \textbf{Reparto (3 primeros)}: 22\,100 actores distintos; solo 1\,114 salen en 10 películas o más, y cubren el 36,4\,\% de
        las apariciones. Si se hace multi-hot de todos, la mayoría de las columnas aparecen una vez: el vector no generaliza.
        Dos soluciones: frecuencia mínima (los demás van a \texttt{<OTRO>}) o \emph{feature hashing} a un tamaño fijo.
  \item \textbf{Palabras clave}: aportan matiz temático («heist», «time travel») pero se solapan con la sinopsis; se pueden añadir
        al texto antes de pasarlo por MiniLM (sin columnas nuevas).
\end{itemize}

\subsection{Año, duración y popularidad}
\begin{itemize}
  \item \textbf{Año y duración}: baratos y con buena cobertura (99,7\,\%; 84,7\,\%, con mediana de 100\,min y p5--p95 de 70--142). Hay que
        normalizarlos con la media y desviación de train (o de catálogo: son metadatos), o agruparlos por décadas.
  \item \textbf{Popularidad} (nº de ratings, \texttt{vote\_count}\dots): el 10\,\% de películas más vistas reúne el 78,8\,\% de los
        ratings y allí la tasa de positivos es más alta (53,3\,\% de $r\ge4$ en el primer decil frente al 30--40\,\% del resto).
        La popularidad ya entra en el entrenamiento por las frecuencias de las películas positivas; con logQ, el modelo estima
        la verdadera $p(i\mid u)$ incluyendo la popularidad. Añadirla como feature da al modelo un atajo para acertar
        recomendando siempre lo mismo, justo el problema que E1 dice combatir.
  \item Una analogía: YouTube mete la «edad del ejemplo» como feature al entrenar y la pone a cero al servir, para que el modelo
        no aprenda un sesgo del pasado \cite{covington2016}. La popularidad podría tratarse igual (un sesgo $b_i$ solo en
        entrenamiento); no es lo mismo que logQ, pero comparte la idea.
  \item Correlación de Spearman entre nº de ratings y año: $-0{,}14$. Las películas de los años 90 tienen la mediana más alta (370 ratings);
        las de los 2010 y 2020, 70 y 56.
\end{itemize}

\subsection{Embedding de \texttt{movieId} frente a solo contenido}
Un embedding por película (tabla $23\,335\times128$, casi 3 millones de parámetros) deja que el modelo memorice qué gusta
a quién. Pero con $1\,491\,641$ positivos de train en la submuestra hay $\approx0{,}50$ positivos por parámetro, y 12\,872 películas
tienen menos de 5 positivos (4\,354, ninguno): su embedding casi no se entrena y su representación sería ruido. La torre de ítem de
YouTube combina identificador de vídeo y características de contenido \cite{yi2019}. Alternativas:
\begin{itemize}
  \item \textbf{Solo contenido}: funciona con películas nuevas, pero pierde la señal colaborativa fina.
  \item \textbf{Contenido + ID}: lo mejor para las populares, pero para una película sin ID entrenado el vector depende de un embedding aleatorio.
  \item \textbf{Contenido + ID con \emph{dropout} del ID}: se apaga el ID al azar durante el entrenamiento para que el contenido tenga
        que sostener el vector; al servir, una película nueva usa un ID «desconocido».
\end{itemize}
No hay datos que decidan esto sin entrenar: es un experimento (con la prueba de D8).

\subsection{Cómo se combinan, y normalizar}
El vector de ítem sale de una red densa sobre la concatenación de [sinopsis (384), géneros (multi-hot), embeddings de director y
actores, año, duración, ¿ID?]. Yi et al. encuentran que \textbf{normalizar con L2} los vectores de usuario y de ítem y añadir una temperatura
mejora la entrenabilidad y la calidad \cite{yi2019}. Para nosotros hay una razón extra: el MMR usa la similitud $\text{Sim}(\mathbf v_i,\mathbf v_j)$
(coseno) y HNSW busca por producto escalar: con vectores normalizados las dos cosas coinciden.

\subsection{¿Congelado o ajustado?}
MiniLM puede quedarse congelado (se precalcula, se cachea y no entra en el bucle de entrenamiento) o ajustarse (\emph{fine-tuning}). Ajustarlo
obliga a recalcular los embeddings cada vez que se actualiza el modelo y complica leer los \texttt{.npy} desde C++ en inferencia. Un compromiso:
congelado y una capa densa entrenable por encima.

% =====================================================================
\section{Las preguntas}
% =====================================================================
\rec\ = mi recomendación, no una decisión. Las preguntas D1--D3 y E1--E4 vienen de \texttt{docs/\allowbreak planning\_\allowbreak bloques\_\allowbreak D\_H.md}
(15/10/2026), con su numeración y su texto; las demás las añado yo. «Afecta C» = decisiones de C relacionadas.

\newcommand{\preg}[3]{\subsection*{#1 \normalfont\small\hfill #3}\textit{#2}\par}
\newcommand{\afecta}[1]{\par\smallskip\noindent\textbf{Afecta C:} #1\par}

% ---------------------------------------------------------------- D
\subsection*{\large Bloque D: split}

\preg{D1. Tipo de split}{¿Cómo partimos?}{F1, D4, D8}
\begin{enumerate}[label=\alph*)]
  \item Aleatorio por rating. El 98,9\,\% del test es anterior al último rating visible: infla las métricas.
  \item Temporal global (test = después del 22/10/2020 aprox.). El 86,3\,\% del test es de usuarios sin historia; solo 439 usuarios elegibles.
  \item Último de cada usuario (\emph{leave-last-out}). Estándar en NCF/SASRec. En el 24,4\,\% de usuarios el corte cae en un empate.
  \item \rec\ Por usuarios: 80/10/10 de usuarios, 80\,\% de su historia como contexto y 20\,\% a predecir. Es el escenario Letterboxd.
\end{enumerate}
\textbf{Modelo}: (d) exige que la torre de usuario lea el historial (F1 b, no un \texttt{userId}). \textbf{Pérdida}: solo los ratings de usuarios
de train generan ejemplos; el contexto de val/test no. \textbf{Evaluación}: solo en (d) todos los usuarios de test son nuevos y hay
2\,051 elegibles. \textbf{E2}: etapa \texttt{03\_split} y tabla \texttt{usuario-grupo} como artefacto.
\afecta{C1 b (los positivos de test dependen del split), C5 a (los ejemplos son solo de los usuarios de train), C8 b, C14 c (la máscara se construye con el historial visible).}

\preg{D2. Orden dentro del usuario}{En los usuarios de test, ¿el 80/20 es cronológico o aleatorio?}{D5, C9}
\begin{enumerate}[label=\alph*)]
  \item \rec\ Cronológico (primer 80\,\% como contexto, último 20\,\% a predecir), con desempate fijo por \texttt{movieId}.
  \item Aleatorio dentro de cada usuario.
  \item Cronológico y, como control, repetir con (b).
\end{enumerate}
\textbf{Evaluación}: en el split por usuarios, el 62,2\,\% de los ratings a predecir se hizo el mismo día que el último de contexto; el
62,5\,\% de todos los ratings es del primer día. \textbf{Pérdida}: ninguna. \textbf{E2}: parámetro \texttt{--orden \{cronologico,aleatorio\}}.
(c) cuesta una segunda pasada; solo se justifica si la diferencia es grande.
\afecta{C9 c (el peso 0,75 al primer día es de entrenamiento; la evaluación va sin pesos), C12 a.}

\preg{D3. Estadísticas solo con train}{¿Qué se calcula solo con los datos de entrenamiento?}{C6, C13, E2, E4}
\begin{enumerate}[label=\alph*)]
  \item \rec\ Todo lo «aprendido de los datos» se ajusta con train (patrón \emph{fit/transform}); para el modelo final, se vuelve a ajustar con train+val.
  \item Siempre solo con train, también en el modelo final.
  \item Con todo: no.
\end{enumerate}
\textbf{Qué entra}: $\bar r_u$ (C6), $P(i)$ (C13), máscara (C14), vocabulario y frecuencia mínima de directores/actores (E2),
media y desviación de año y duración (E4). \textbf{Fuga medida} (si se hiciera con todo): 1,46\,\% de etiquetas de train cambian y
1\,730 películas reciben un $P(i)$ gracias a test. \textbf{E2}: cada transformador guarda sus parámetros ajustados en un artefacto.
\afecta{C6 a, C13 b, C14 c.}

\preg{D4. Validación}{¿Con qué se eligen hiperparámetros y se para el entrenamiento?}{D1, D7}
\begin{enumerate}[label=\alph*)]
  \item \rec\ Usuarios de validación apartados (como el test: 80\,\% contexto, 20\,\% a predecir), 10\,\% de los usuarios.
  \item Un 10\,\% de los ratings de train, al azar.
  \item Validación cruzada por usuarios (k carpetas).
  \item Validar solo con una submuestra de los usuarios de validación (más rápido).
\end{enumerate}
\textbf{Evaluación}: criterio de parada, Recall@200 o NDCG@20 sobre el catálogo completo (47 millones de productos escalares por pasada en la
submuestra, asumible). (b) mezcla el futuro en el entrenamiento y no valida «usuario nuevo». (c) multiplica el coste por $k$.
\afecta{C7 c (la validación también tiene las dos definiciones de relevante), C8 b.}

\preg{D5. Primer día y empates}{¿Qué hacemos con el 59\,\% de usuarios que lo valoran todo el primer día?}{C9, D2}
\begin{enumerate}[label=\alph*)]
  \item \rec\ Nada especial en el split: aceptar el orden y desempatar con \texttt{movieId}; documentar que el orden dentro del primer día no es información.
  \item Que el objetivo sean solo ratings posteriores a las primeras 24\,h (solo usuarios con varias sesiones).
  \item Para los usuarios de «un solo día», 80/20 aleatorio; para el resto, cronológico.
  \item Quitar a los usuarios cuyo corte cae en un empate de segundo (24,4\,\%).
\end{enumerate}
\textbf{Evaluación}: (b) deja un test que mide otra cosa (ya no es Letterboxd) y menos usuarios; (d) tira casi un cuarto de los usuarios. (c)
es defendible pero añade una regla que explicar. \textbf{Pérdida}: ya se trata con C9 c.
\afecta{C9 c, C12 a.}

\preg{D6. Métricas y $K$}{¿Qué se mide y con qué $K$?}{C7, F, MMR}
\begin{enumerate}[label=\alph*)]
  \item Recall@20 y NDCG@20.
  \item Recall@200 y NDCG@20.
  \item \rec\ $K\in\{10,20,50,200\}$ más cobertura@20 y novedad@20, siempre sobre el catálogo completo; principal: Recall@200 (retrieval) y NDCG@20 (lista final, antes y después del MMR).
  \item Métricas con 100 negativos muestreados.
\end{enumerate}
\textbf{Evaluación}: techo de Recall@10 en el split por usuarios: 77,1\,\%; @20: 89,9\,\%; @200: 100\,\%. (d) no es comparable entre modelos
(guía C). \textbf{Pérdida}: ninguna, pero las métricas se calculan con las dos definiciones de C7 c. \textbf{E2}: intervalos de confianza por \emph{bootstrap} de usuarios.
\afecta{C7 c, C8 b (los usuarios sin positivos salen), C14 c (en inferencia se excluye lo ya valorado).}

\preg{D7. Tamaño de las particiones}{¿Qué fracción de usuarios para val y test?}{D1, D9}
\begin{enumerate}[label=\alph*)]
  \item \rec\ 80/10/10 de usuarios (en la submuestra, $\approx$2\,000 usuarios en val y en test; en el completo, $\approx$20\,000).
  \item 90/5/5.
  \item Un número fijo de usuarios (p.\,ej. 10\,000 como en Mult-VAE \cite{liang2018}) en el conjunto completo.
\end{enumerate}
\textbf{Evaluación}: con 2\,051 elegibles y 35\,019 positivos hay ruido (la semilla mueve $\pm38$ usuarios elegibles); más usuarios de test = menos ruido, a costa de
menos entrenamiento. \textbf{Modelo}: con (a) hay 15\,960 usuarios de entrenamiento en la submuestra.
\afecta{C5 a, C8 b.}

\preg{D8. Arranque en frío}{¿Qué hacemos con películas y usuarios de test sin historia en train?}{E6, D1}
\begin{enumerate}[label=\alph*)]
  \item Ignorarlo (dejarlos en el test tal cual).
  \item Quitar del test las películas sin train (solo el 0,15\,\% en el split por usuarios).
  \item Reportar aparte: métricas globales y métricas sobre ratings en películas con $<5$ (o $<20$) en train.
  \item \rec\ c y, además, un experimento de películas retiradas: quitar todos los ratings de train de un 5\,\% de películas y medir sobre ellas.
\end{enumerate}
\textbf{Evaluación}: con el split por usuarios solo el 1,54\,\% del test es de películas con $<5$ ratings en train: sin el experimento d, la promesa
de arranque en frío de E1 no se puede probar. \textbf{Usuarios}: Letterboxd ya es el escenario del split (d) de D1. \textbf{Modelo}: la prueba decide E6.
\afecta{C13 b (las películas sin positivos tienen $P(i)$ solo por el $+1$), C11 b.}

\preg{D9. Submuestra o completo}{¿Con qué conjunto se desarrolla y con cuál se informa?}{D7, D10}
\begin{enumerate}[label=\alph*)]
  \item Submuestra para todo, completo solo al final.
  \item \rec\ Submuestra para desarrollar y, antes de cada entrega, repetir las 2--3 configuraciones finales en el completo.
  \item Completo desde el principio.
\end{enumerate}
\textbf{Evaluación}: el frío de la submuestra es peor que el real (15\,422 de 23\,335 películas con $<20$ ratings en train); los cortes
temporales coinciden. \textbf{E2}: parámetro \texttt{--muestra}.
\afecta{C13 b, C8 b.}

\preg{D10. Semillas y reproducibilidad}{¿Cómo hacemos el split reproducible?}{D7, H7}
\begin{enumerate}[label=\alph*)]
  \item Una semilla fija (42).
  \item \rec\ Asignar cada usuario a su grupo con una función hash estable de \texttt{userId} y la semilla; guardar la tabla; repetir los resultados finales con 5 semillas y dar media y desviación.
  \item Semilla distinta para cada componente (split, inicialización, orden de lotes).
\end{enumerate}
\textbf{Evaluación}: con 5 semillas, los ratings de test varían $66\,565\pm1\,569$ (usuarios) y $315\,203\pm223$ (aleatorio). \textbf{E2}: la tabla de grupos y
todos los umbrales van en el archivo de configuración.
\afecta{C1 b, C13 b (todo lo que depende de train cambia con la semilla).}

% ---------------------------------------------------------------- E
\subsection*{\large Bloque E: features de la torre de película}

\preg{E1. Embedding de la sinopsis}{¿Cómo se convierte el texto en vector?}{E8}
\begin{enumerate}[label=\alph*)]
  \item \rec\ MiniLM (\texttt{all-MiniLM-L6-v2}, $d=384$), precalculado.
  \item TF-IDF + SVD.
  \item Ambos (ablación).
\end{enumerate}
\textbf{Modelo}: la sinopsis ocupa 384 de las entradas de la torre; en \texttt{float16} son 17,9\,MB para 23\,335 películas, en \texttt{float32}, 35,8\,MB. Mediana de 46
palabras (cabe en la ventana). \textbf{Pérdida/evaluación}: ninguna directa; (c) da una comparación gratis. \textbf{E2}: caché \texttt{.npy}. 127 sinopsis vacías: vector de ceros y bandera.
\afecta{C13 b (el vector solo describe contenido: no usa ratings).}

\preg{E2. Créditos}{¿Qué parte de director y reparto?}{E9, D3}
\begin{enumerate}[label=\alph*)]
  \item Nada.
  \item \rec\ Director + 3--5 primeros actores, con frecuencia mínima (p.\,ej. $\ge5$ películas; el resto a \texttt{<OTRO>}) o \emph{hashing}.
  \item Todo el reparto y el equipo.
\end{enumerate}
\textbf{Modelo}: 7\,827 directores y 22\,100 actores (solo 1\,114 con $\ge10$ películas; 36,4\,\% de las apariciones); sin frecuencia mínima, columnas que casi nunca se entrenan. Cobertura: 84,8\,\% de
películas (97,0\,\% de ratings), pero 9,2\,\% de las posteriores a 2017. \textbf{D3}: la frecuencia mínima y el vocabulario se cuentan con las películas de train.
\afecta{C13 b (las películas con $P(i)=0$ dependerán de estas features), C11 b.}

\preg{E3. Palabras clave}{¿Se incluyen?}{E1}
\begin{enumerate}[label=\alph*)]
  \item No.
  \item \rec\ Añadirlas al texto antes de MiniLM (sin features nuevas).
  \item Multi-hot con frecuencia mínima.
\end{enumerate}
\textbf{Cobertura}: 84,7\,\% de las películas (98,5\,\% de los ratings); el 91,9\,\% de las posteriores a 2017 (más que los créditos, porque la API de TMDB sí las devuelve).
\textbf{Modelo}: (b) no cambia la arquitectura; (c) añade otro vocabulario con la misma pega que E2.
\afecta{---}

\preg{E4. Features numéricas}{¿Año, duración, popularidad, \texttt{vote\_average}?}{D3, C13}
\begin{enumerate}[label=\alph*)]
  \item \rec\ Año y duración (normalizados con train o por décadas), sin popularidad ni \texttt{vote\_average}.
  \item Todas.
  \item (a) más la popularidad solo como sesgo de entrenamiento.
  \item Ninguna.
\end{enumerate}
\textbf{Modelo}: año 99,7\,\% de cobertura; duración 84,7\,\% (8,9\,\% en las $\ge2017$). La popularidad empuja al modelo a recomendar lo popular: el 10\,\% más visto reúne el 78,8\,\% de los ratings, y
con C13 b logQ ya trae esa información a la pérdida. \texttt{vote\_average} viene de otros usuarios de TMDB y mezcla calidad con popularidad.
\afecta{C13 b (logQ y popularidad como feature se pisan), C1 b.}

\preg{E5. Géneros (decisión B3)}{¿MovieLens o TMDB?}{E9, inferencia}
\begin{enumerate}[label=\alph*)]
  \item Solo MovieLens (19 géneros; 99,0\,\% de cobertura).
  \item Solo TMDB (23 géneros; 98,7\,\%).
  \item \rec\ Unión: TMDB primero y MovieLens como relleno (hay 258 películas con géneros solo en MovieLens), quitando IMAX (formato) y Foreign (no es un género); vocabulario común de 16 mapeado.
  \item Solo los 16 géneros comunes.
\end{enumerate}
\textbf{Modelo}: Jaccard medio 0,75; listas idénticas en el 50,2\,\%. Las películas nuevas y las de Letterboxd solo tendrán géneros de TMDB: entrenar con géneros de MovieLens crea una discrepancia entre entrenamiento y servicio.
38 películas sin géneros en ninguna: vector de ceros y bandera.
\afecta{C3 c (las anti-preferencias de la torre de usuario usan los géneros de las películas), C10 a.}

\preg{E6. Identificador de película}{¿Embedding de \texttt{movieId} o solo contenido?}{D8, F1}
\begin{enumerate}[label=\alph*)]
  \item Solo contenido.
  \item Contenido + embedding de ID.
  \item \rec\ Contenido + ID con \emph{dropout} del ID (y un ID «desconocido» para películas nuevas).
\end{enumerate}
\textbf{Modelo}: tabla de 2\,986\,880 parámetros frente a 1\,491\,641 positivos de train; 12\,872 películas con $<5$ positivos y 4\,354 con ninguno. \textbf{Evaluación}: se decide con el experimento de D8 d; una
película nueva sin ID entrenado solo se sostiene con contenido. \textbf{Pérdida}: logQ no cambia, pero el ID permite que el modelo memorice la popularidad.
\afecta{C13 b, C5 a (sin tope por usuario, el ID aprende sobre todo de los muy activos).}

\preg{E7. Combinación y normalización}{¿Cómo se juntan las features y se normalizan los vectores?}{E1, MMR}
\begin{enumerate}[label=\alph*)]
  \item Concatenar y una red densa; sin normalizar $\mathbf v_i$.
  \item \rec\ Concatenar y red densa; normalizar con L2 $\mathbf u$ y $\mathbf v_i$ y usar una temperatura $T$ (\cite{yi2019}).
  \item Normalizar solo $\mathbf v_i$.
\end{enumerate}
\textbf{Modelo y pérdida}: la temperatura es un hiperparámetro de la softmax (D4). \textbf{Inferencia}: con vectores normalizados, el producto escalar de HNSW y el coseno del MMR coinciden.
\afecta{C1 b, C13 b (se resta $\log P(i)$ a la puntuación con $T$).}

\preg{E8. Encoder congelado o ajustado}{¿Se entrena MiniLM?}{E1}
\begin{enumerate}[label=\alph*)]
  \item \rec\ Congelado, con una capa densa entrenable encima.
  \item Ajustar la última(s) capa(s).
  \item Ajustar todo.
\end{enumerate}
\textbf{Coste}: (a) no entra en el bucle de entrenamiento y se lee con \texttt{.npy} desde C++; (b)/(c) obligan a recalcular los 23\,335 vectores tras cada cambio.
\afecta{---}

\preg{E9. Datos que faltan}{¿Qué hacemos con las películas sin director, reparto o duración (posteriores a 2017)?}{E2, E4}
\begin{enumerate}[label=\alph*)]
  \item \rec\ Token «desconocido» y una bandera \texttt{falta\_credito}; con \emph{dropout} de los créditos al entrenar (para que el modelo vea películas sin ellos).
  \item Descargar los créditos que faltan con la API de TMDB (como se hizo con las sinopsis).
  \item Excluir esas películas del entrenamiento.
\end{enumerate}
\textbf{Modelo}: solo el 9,2\,\% de las películas de 2017 en adelante tienen director; sin tratamiento, la torre aprende «sin director = reciente». \textbf{Datos}: (b) cuesta una descarga
(3\,361 películas) pero cambia la cobertura; se puede hacer después. \textbf{Evaluación}: se mira en el experimento de D8 d.
\afecta{C11 b, C13 b.}

% ---------------------------------------------------------------- dependencias
\subsection*{Resumen de dependencias}
\begin{center}
\small
\begin{tabularx}{\linewidth}{lYY}
\toprule
Pregunta & Depende de\dots & Porque\dots \\
\midrule
D1 & F1 (y C1) & el split por usuarios exige torre de usuario basada en historial \\
D2 & C9 & el orden en el primer día es el de tecleo \\
D3 & C6, C13, C14 & definen qué se calcula con train \\
D4, D6 & C7, C8 & las dos definiciones de relevante y los usuarios elegibles \\
D5 & C9 & el 59\,\% de usuarios en un solo día \\
D8 & E6, D1 & el frío de películas solo se mide con películas retiradas \\
E1, E2, E3 & D3 & vocabularios, frecuencias mínimas, idf: solo con train \\
E4 & C13 & logQ y popularidad como feature se pisan \\
E5 & B3 & fija la decisión B3 \\
E6 & D8, F1 & decide si la torre funciona con solo contenido \\
E9 & E2, E4 & las películas posteriores a 2017 carecen de créditos \\
\bottomrule
\end{tabularx}
\end{center}
Orden propuesto: (1) D1, F1 y E6 a la vez (qué sabe cada torre y con qué se mide); (2) D3 y E1--E4 (qué se ajusta y con qué datos);
(3) D2, D4--D5 y D7 (detalles del split y empates); (4) E5, E7--E9 (géneros, normalización y datos que faltan); (5) D6, D8--D10 al escribir
el pipeline de evaluación y el archivo de configuración.

% =====================================================================
\section*{Glosario}
\addcontentsline{toc}{section}{Glosario}
% =====================================================================
\begin{description}[leftmargin=3.2cm, style=nextline, itemsep=1pt]
  \item[Arranque en frío] Un usuario o una película sin historia en train.
  \item[Contexto / \emph{fold-in}] Parte del historial de un usuario de test que se da al modelo para construir su vector; no se predice.
  \item[Elegible] Usuario de test con $\ge1$ positivo en test y $\ge1$ en su historial visible.
  \item[\emph{Fit/transform}] Ajustar una transformación con train y aplicar el resultado ya fijado al resto.
  \item[Fuga (\emph{leakage})] Usar en entrenamiento información que solo existe en test.
  \item[Generalización fuerte] Los usuarios de test no están en el entrenamiento.
  \item[\emph{Hashing}] Reducir un vocabulario enorme a un tamaño fijo con una función hash.
  \item[\emph{Leave-last-out}] Último rating de cada usuario a test.
  \item[MiniLM] Red preentrenada que convierte un texto en un vector de 384 números.
  \item[Techo de Recall@$K$] Máximo Recall@$K$ alcanzable: media de $\min(K,|R_u|)/|R_u|$.
  \item[TF-IDF / SVD] Ponderación de palabras por su rareza y compresión por valores singulares.
\end{description}

\begin{thebibliography}{99}
\bibitem{carbonell1998} J.~Carbonell, J.~Goldstein. \emph{The Use of MMR, Diversity-Based Reranking for Reordering Documents and Producing Summaries}. SIGIR, 1998.
\bibitem{covington2016} P.~Covington, J.~Adams, E.~Sargin. \emph{Deep Neural Networks for YouTube Recommendations}. RecSys, 2016.
\bibitem{harper2015} F.~M.~Harper, J.~A.~Konstan. \emph{The MovieLens Datasets: History and Context}. ACM TiiS, 2015.
\bibitem{he2017} X.~He, L.~Liao, H.~Zhang, L.~Nie, X.~Hu, T.-S.~Chua. \emph{Neural Collaborative Filtering}. WWW, 2017.
\bibitem{kang2018} W.-C.~Kang, J.~McAuley. \emph{Self-Attentive Sequential Recommendation}. ICDM, 2018.
\bibitem{liang2018} D.~Liang, R.~G.~Krishnan, M.~D.~Hoffman, T.~Jebara. \emph{Variational Autoencoders for Collaborative Filtering}. WWW, 2018.
\bibitem{steck2019} H.~Steck. \emph{Embarrassingly Shallow Autoencoders for Sparse Data}. WWW, 2019.
\bibitem{yi2019} X.~Yi, J.~Yang, L.~Hong, D.~Z.~Cheng, L.~Heldt, A.~Kumthekar, Z.~Zhao, L.~Wei, E.~Chi. \emph{Sampling-Bias-Corrected Neural Modeling for Large Corpus Item Recommendations}. RecSys, 2019.
\end{thebibliography}

\end{document}
